\section{Why the defender's word isn't enough}
\label{sec:5.2}
The exception for defense is where wars of choice live. Every belligerent claims to be the defender, so the claim is the belligerent's own, and it is discounted like any interested claim. Iraq is the instance: the judgment that Iraq had mobile biological-weapons plants, the core of its claim to biological weapons, rested almost exclusively on a single informant, known as Curveball, whom no American intelligence officer had met before the war: a corroboration with one origin \autocite{wmdcommission2005}. The exception therefore runs in two tiers. The first tier, force against an attack under way on the defender's own account, is the adoptable term's Article 51 as courts read it; what the corroborated-defense departure adds is the second tier, under which the licence persists only on disinterested corroboration.

That is the first departure, and the one the rest of the chapter leans on. The floor lifts its prohibition on force for self-defense on the defender's own claim only while an attack is under way, and keeps it lifted only on disinterested corroboration. The Charter lets a state judge its own claim until the Security Council acts.

\textbf{The first tier acts at once:} the force repels an armed attack under way against the defender's territory, forces or population, including attacks on the forces, launchers and staging areas conducting it, and assistance given at the defender's request. It ends when the attack ends. It covers the least force that repels the attack: the necessity and proportionality the Court has read into Article 51 since \emph{Nicaragua} \autocite{icj1986nicaragua,icj2003oilplatforms}. It does not cover a campaign to destroy the attacker, to occupy its territory, or to punish the population it came from. The first waits on the second tier; the second and third fall under no exception. An invasion can't wait for a vote. It acts on the defender's own account, because in the hours that matter there is no other, and that is the tier's weakness: Gleiwitz in 1939, the Mukden railway in 1931 and the Gulf of Tonkin in 1964 were each an attack under way on the defender's account, and each opened a war of choice. The first tier reads administration, not title, because deciding whose territory it is would be an administered judgment and the first tier doesn't need one. Administration is a line a disinterested body has already drawn. In 2005 the Eritrea–Ethiopia Claims Commission held that Eritrea broke Article 2(4) by attacking territory then under Ethiopia's peaceful administration, because self-defense “cannot be invoked to settle territorial disputes” \autocite{eecc2005jusadbellum}, though the boundary delimitation had put Badme, the town at issue, on Eritrea's side \autocite{bbc2003badme}. So the rule is the Charter as a tribunal has read it, not a departure. Territory a government has administered in peace counts as its own for the first tier, whoever claims it, unless it took that territory by force in the conflict at issue or a disinterested body has found its hold an occupation. Who holds the elections, runs the courts and issues the papers is a fact on the record. An attacker's claim of title is an interested claim and is discounted as one; its remedy is the second tier, never force. So the first tier protects any authority holding ground no disinterested body has found occupied, including one a patron keeps in place; a request from such an authority for its patron's help is the patron's own case and counts for nothing, and where the patron's forces are already there, I have no test. Where the line of administration is itself contested on the record, the first tier doesn't apply across it and the second governs from the start, and even there the exclusion reaches an attack on the territory and not an attack on the population living there, which the first tier covers wherever it occurs.

The first tier ends when the attack ends, and whether it has ended is a fact the defender controls the evidence about. The burden is therefore the defender's. Its uncorroborated claim that the attack continues moves nothing; the licence persists on disinterested corroboration, or lapses.

A deployment can do three things about a manufactured incident. It holds the record of what the defender claimed and when, so a licence obtained on a staged incident is a licence on a record. It takes evidence from channels the defender doesn't compose, which is a condition of licensing. And it lets the licence lapse when outside evidence arrives, which is the second tier's job. A refuser that remembers having been misled and raises its discount on the source does all three faster.

\textbf{The second tier confirms or lifts the first later, and for the genocide and occupation exceptions originates a licence of its own.} What it needs is a disinterested finding, of one of four kinds: a judgment or advisory opinion of the International Court of Justice, or a judgment of a regional court with adversarial process; the finding of a fact-finding mission or commission of inquiry mandated by a body no party to the conflict controls; a determination of the Security Council; or a resolution adopted by two-thirds of the General Assembly members present and voting, where the facts it rests on aren't in dispute. A Council authorization counts because it needs the permanent members not to block it. The Council's silence on a defense claim counts for nothing, because a veto is a stake. The four kinds are not equal, and Appendix~\ref{sec:G.1} says how they are weighed.

Neither tier is a certificate issued once. Whether force still falls under an exception is reassessed as evidence arrives, by the system and by whoever holds the record. What the licence records is that ongoing judgment: which exception the force rests on now, on what evidence, what would end it, and each point at which the justification weakened and the force continued.
